% Notes in bubble
% put 'final' instead of 'draft' to inhibite the package
% put 'disable' to inhibite the package
\usepackage[colorinlistoftodos]{todonotes}%textwidth=35mm
%\usepackage[textwidth=35mm,colorinlistoftodos,disable]{todonotes}
%\setlength{\marginparwidth}{1in}
%\setlength{\marginparwidth}{2cm}
\definecolor{antiquewhite}{rgb}{0.98, 0.92, 0.84}
\definecolor{buff}{rgb}{0.94, 0.86, 0.51}
\definecolor{palecopper}{rgb}{0.85, 0.54, 0.4}
\definecolor{fluorescentyellow}{rgb}{0.8, 1.0, 0.0}
\definecolor{bole}{rgb}{0.47, 0.27, 0.23}
% A page-wide note with stuff to be added later.
\newcommand{\towrite}[1]{\leavevmode\todo[inline,color=antiquewhite]{\textbf{TO WRITE: }#1}}

% Notes in a coloured bubble.
\newcommand{\tonote}[1]{\todo[size=\tiny,color=buff]{\textbf{NOTE: }  #1}}
\newcommand{\tofix}[1]{\todo[size=\tiny,color=palecopper]{\textbf{TO FIX: } #1}}
\newcommand{\toask}[1]{\todo[size=\tiny,color=fluorescentyellow]{\textbf{QUESTION: } #1}} 

\usepackage{amsmath,float}
\usetikzlibrary{decorations.markings}








\usepackage[vcentermath]{youngtab}
\input xy
\xyoption{all}


\usepackage{pgfkeys}
\usepackage{pgfopts}
\usepackage{ytableau}

\definecolor{cornellred}{rgb}{0.7, 0.11, 0.11}
\definecolor{britishracinggreen}{rgb}{0.0, 0.26, 0.15}
\definecolor{cobalt}{rgb}{0.0, 0.28, 0.67}
\DeclareSymbolFont{usualmathcal}{OMS}{cmsy}{m}{n}
\DeclareSymbolFontAlphabet{\mathcal}{usualmathcal}

\newcommand{\TT}{\mathbf{T}}
\newcommand{\bfk}{\mathbf{k}}


%\lstset{language = Macaulay2}

\newcommand{\BA}{{\mathbb{A}}}
\newcommand{\BB}{{\mathbb{B}}}
\newcommand{\BC}{{\mathbb{C}}}
\newcommand{\BD}{{\mathbb{D}}}
\newcommand{\BE}{{\mathbb{E}}}
\newcommand{\BF}{{\mathbb{F}}}
\newcommand{\BG}{{\mathbb{G}}}
\newcommand{\BH}{{\mathbb{H}}}
\newcommand{\BI}{{\mathbb{I}}}
\newcommand{\BJ}{{\mathbb{J}}}
\newcommand{\BK}{{\mathbb{K}}}
\newcommand{\BL}{{\mathbb{L}}}
\newcommand{\BM}{{\mathbb{M}}}
\newcommand{\BN}{{\mathbb{N}}}
\newcommand{\BO}{{\mathbb{O}}}
\newcommand{\BP}{{\mathbb{P}}}
\newcommand{\BQ}{{\mathbb{Q}}}
\newcommand{\BR}{{\mathbb{R}}}
\newcommand{\BS}{{\mathbb{S}}}
\newcommand{\BT}{{\mathbb{T}}}
\newcommand{\BU}{{\mathbb{U}}}
\newcommand{\BV}{{\mathbb{V}}}
\newcommand{\BW}{{\mathbb{W}}}
\newcommand{\BX}{{\mathbb{X}}}
\newcommand{\BY}{{\mathbb{Y}}}
\newcommand{\BZ}{{\mathbb{Z}}}

\newcommand{\CA}{{\mathcal A}}
\newcommand{\CB}{{\mathcal B}}
\newcommand{\CC}{{\mathcal C}}
\newcommand{\CD}{{\mathcal D}}
\newcommand{\CE}{{\mathcal E}}
\newcommand{\CF}{{\mathcal F}}
\newcommand{\CG}{{\mathcal G}}
\newcommand{\CH}{{\mathcal H}}
\newcommand{\CI}{{\mathcal I}}
\newcommand{\CJ}{{\mathcal J}}
\newcommand{\CK}{{\mathcal K}}
\newcommand{\CL}{{\mathcal L}}
\newcommand{\CM}{{\mathcal M}}
\newcommand{\CN}{{\mathcal N}}
\newcommand{\CO}{{\mathcal O}}
\newcommand{\CP}{{\mathcal P}}
\newcommand{\CQ}{{\mathcal Q}}
\newcommand{\CR}{{\mathcal R}}
\newcommand{\CS}{{\mathcal S}}
\newcommand{\CT}{{\mathcal T}}
\newcommand{\CU}{{\mathcal U}}
\newcommand{\CV}{{\mathcal V}}
\newcommand{\CW}{{\mathcal W}}
\newcommand{\CX}{{\mathcal X}}
\newcommand{\CY}{{\mathcal Y}}
\newcommand{\CZ}{{\mathcal Z}}

\newcommand{\Fa}{{\mathfrak{a}}}
\newcommand{\Fb}{{\mathfrak{b}}}
\newcommand{\Fc}{{\mathfrak{c}}}
\newcommand{\Fd}{{\mathfrak{d}}}
\newcommand{\Fe}{{\mathfrak{e}}}
\newcommand{\Ff}{{\mathfrak{f}}}
\newcommand{\Fg}{{\mathfrak{g}}}
\newcommand{\Fh}{{\mathfrak{h}}}
\newcommand{\Fi}{{\mathfrak{i}}}
\newcommand{\Fj}{{\mathfrak{j}}}
\newcommand{\Fk}{{\mathfrak{k}}}
\newcommand{\Fl}{{\mathfrak{l}}}
\newcommand{\Fm}{{\mathfrak{m}}}
\newcommand{\Fn}{{\mathfrak{n}}}
\newcommand{\Fo}{{\mathfrak{o}}}
\newcommand{\Fp}{{\mathfrak{p}}}
\newcommand{\Fq}{{\mathfrak{q}}}
\newcommand{\Fr}{{\mathfrak{r}}}
\newcommand{\Fs}{{\mathfrak{s}}}
\newcommand{\Ft}{{\mathfrak{t}}}
\newcommand{\Fu}{{\mathfrak{u}}}
\newcommand{\Fv}{{\mathfrak{v}}}
\newcommand{\Fw}{{\mathfrak{w}}}
\newcommand{\Fx}{{\mathfrak{x}}}
\newcommand{\Fy}{{\mathfrak{y}}}
\newcommand{\Fz}{{\mathfrak{z}}}

\newcommand{\FA}{{\mathfrak{A}}}
\newcommand{\FB}{{\mathfrak{B}}}
\newcommand{\FC}{{\mathfrak{C}}}
\newcommand{\FD}{{\mathfrak{D}}}
\newcommand{\FE}{{\mathfrak{E}}}
\newcommand{\FF}{{\mathfrak{F}}}
\newcommand{\FG}{{\mathfrak{G}}}
\newcommand{\FH}{{\mathfrak{H}}}
\newcommand{\FI}{{\mathfrak{I}}}
\newcommand{\FJ}{{\mathfrak{J}}}
\newcommand{\FK}{{\mathfrak{K}}}
\newcommand{\FL}{{\mathfrak{L}}}
\newcommand{\FM}{{\mathfrak{M}}}
\newcommand{\FN}{{\mathfrak{N}}}
\newcommand{\FO}{{\mathfrak{O}}}
\newcommand{\FP}{{\mathfrak{P}}}
\newcommand{\FQ}{{\mathfrak{Q}}}
\newcommand{\FR}{{\mathfrak{R}}}
\newcommand{\FS}{{\mathfrak{S}}}
\newcommand{\FT}{{\mathfrak{T}}}
\newcommand{\FU}{{\mathfrak{U}}}
\newcommand{\FV}{{\mathfrak{V}}}
\newcommand{\FW}{{\mathfrak{W}}}
\newcommand{\FX}{{\mathfrak{X}}}
\newcommand{\FY}{{\mathfrak{Y}}}
\newcommand{\FZ}{{\mathfrak{Z}}}

%\newcommand{\red}{\mathrm{red}}
\newcommand{\ddd}{\mathbf{d}}
\newcommand{\ord}{\mathrm{ord}}
\newcommand{\divisor}{\mathrm{div}}
\newcommand{\codim}{\mathrm{codim}}
\newcommand{\mult}{\mathrm{mult}}
\newcommand{\supp}{\mathrm{supp}}
\newcommand{\simto}{\,\widetilde{\to}\,}
\newcommand{\gr}{\mathrm{gr}}

\newcommand{\dq}{{q \frac{d}{dq}}}
\newcommand{\pt}{{\mathsf{pt}}}
\newcommand{\ch}{{\mathrm{ch}}}
\newcommand{\td}{{\mathrm{td}}}
\newcommand{\fix}{\mathsf{fix}}
\newcommand{\mov}{\mathsf{mov}}
\newcommand{\vol}{{\mathrm{vol}}}
\newcommand{\area}{{\mathrm{area}}}

\DeclareMathOperator{\Hilb}{Hilb}
\DeclareMathOperator{\Eu}{Eu}
\DeclareMathOperator{\Sets}{Sets}
\DeclareMathOperator{\Sch}{Sch}
\DeclareMathOperator{\Quot}{Quot}
\DeclareMathOperator{\con}{con}
\DeclareMathOperator{\red}{red}
\DeclareMathOperator{\Der}{Der}
\DeclareMathOperator{\nilp}{nilp}
\DeclareMathOperator{\coker}{coker}
\DeclareMathOperator{\depth}{depth}
\DeclareMathOperator{\ann}{ann}
\DeclareMathOperator{\Ann}{Ann}
\DeclareMathOperator{\AP}{AP}
\DeclareMathOperator{\BPS}{BPS}
\DeclareMathOperator{\Td}{Td}
\DeclareMathOperator{\BPSs}{\CB\CP\CS}
\DeclareMathOperator{\IC}{IC}
\DeclareMathOperator{\ICs}{\CI\CC}
\DeclareMathOperator{\Sym}{Sym}
\DeclareMathOperator{\Mod}{Mod}
\DeclareMathOperator{\mot}{mot}
\DeclareMathOperator{\coh}{coh}
\DeclareMathOperator{\crr}{cr}
\DeclareMathOperator{\Per}{Per}
\DeclareMathOperator{\Coh}{Coh}
\DeclareMathOperator{\QCoh}{QCoh}
\DeclareMathOperator{\MHS}{MHS}
\DeclareMathOperator{\MMHS}{MMHS}
\DeclareMathOperator{\MHM}{MHM}
\DeclareMathOperator{\MMHM}{MMHM}
\DeclareMathOperator{\Chow}{Chow}
%\DeclareMathOperator{\Tr}{Tr}
\DeclareMathOperator{\Frac}{Frac}
\DeclareMathOperator{\HTop}{HTop}
\DeclareMathOperator{\Orb}{Orb}
\DeclareMathOperator{\Span}{Span}
\DeclareMathOperator{\vd}{vd}
\DeclareMathOperator{\ob}{ob}
\DeclareMathOperator{\Ob}{Ob}
\DeclareMathOperator{\Obs}{Obs}
\DeclareMathOperator{\Boxtimes}{\mathlarger{\mathlarger{\boxtimes}}}

\DeclareMathOperator{\Art}{Art}

%%%%%%%%%%%%%%%%%%%%%%%%%% Davison-Ricolfi commands %%%%%%

\DeclareMathOperator{\vir}{\mathrm{vir}}
\DeclareMathOperator{\twist}{\text{sh}}
\DeclareMathOperator{\Exp}{Exp}
\DeclareMathOperator{\Var}{Var}
\DeclareMathOperator{\KK}{K}
\DeclareMathOperator{\Db}{D^b}
\DeclareMathOperator{\D}{D}
\DeclareMathOperator{\Dlf}{D^{\geq,\textrm{lf}}}
%\DeclareMathOperator{\Rep}{rMod}
\DeclareMathOperator{\Rep}{Rep}
\DeclareMathOperator{\GL}{GL}
\DeclareMathOperator{\PGL}{PGL}
\DeclareMathOperator{\Mat}{Mat}
\DeclareMathOperator{\mon}{mon}
%\DeclareMathOperator{\div}{div}
%\DeclareMathOperator{\ord}{ord}
\DeclareMathOperator{\tw}{tw}
\DeclareMathOperator{\cyc}{cyc}
\DeclareMathOperator{\HN}{HN}
\DeclareMathOperator{\rk}{rk}
%\DeclareMathOperator{\op}{op}
%\DeclareMathOperator{\tr}{tr}
\DeclareMathOperator{\St}{St}
\DeclareMathOperator{\length}{length}
\DeclareMathOperator{\colength}{colength}

\DeclareMathOperator{\NCHilb}{NCHilb}
\DeclareMathOperator{\NCQuot}{Quot}
\DeclareMathOperator{\colim}{colim}
\renewcommand{\ss}{\operatorname{ss}}

\newcommand{\phim}[1]{\phi^{\mathrm{mon}}_{#1}}
\newcommand{\muhat}{\hat{\mu}}
\newcommand{\derived}{\mathbf{D}}
\newcommand{\ff}{\mathbf{f}}
\newcommand{\BBS}{\mathrm{BBS}}
\newcommand{\HC}{\mathsf{hc}}
\newcommand{\HO}{\mathrm{H}}
\newcommand{\fra}{\mathrm{fr}}
\newcommand{\crit}{\operatorname{crit}}
\newcommand{\zetass}{\zeta\textrm{-ss}}
\newcommand{\zetast}{\zeta\textrm{-st}}
\newcommand*{\defeq}{\mathrel{\vcenter{\baselineskip0.5ex \lineskiplimit0pt
                     \hbox{\scriptsize.}\hbox{\scriptsize.}}}%
                     =}

%%%%%%%%%%%%%%%%%%%%%%%%%%%%%%%%%%%%%%%%%%%%%%%%%%%%%%%%%%%%%

\newcommand{\into}{\hookrightarrow}
\newcommand{\onto}{\twoheadrightarrow}

\DeclareFontFamily{OT1}{rsfs}{}
\DeclareFontShape{OT1}{rsfs}{n}{it}{<-> rsfs10}{}
\DeclareMathAlphabet{\curly}{OT1}{rsfs}{n}{it}
\renewcommand\hom{\mathscr{H}\kern-0.3em\mathit{om}}
\newcommand\ext{\curly Ext}
\newcommand\Ext{\operatorname{Ext}}
\newcommand\Hom{\operatorname{Hom}}
\newcommand\End{\operatorname{End}}

\DeclareMathOperator{\lHom}{\mathscr{H}\kern-0.3em\mathit{om}}
\DeclareMathOperator{\RRlHom}{\mathbf{R}\kern-0.025em\mathscr{H}\kern-0.3em\mathit{om}}
\DeclareMathOperator{\RHom}{{\mathbf{R}\mathrm{Hom}}}
\DeclareMathOperator{\lExt}{{\mathscr{E}\kern-0.2em\mathit{xt}}}
\DeclareMathOperator{\Def}{Def}
\DeclareMathOperator{\Jac}{Jac}
\DeclareMathOperator{\pr}{pr}
\DeclareMathOperator{\At}{At}

\newcommand{\p}{\mathbb{P}}
\newcommand{\RR}{\mathbf R}
\newcommand{\hint}{\textbf{Hint}}
\newcommand\Id{\operatorname{Id}}
\newcommand\Spec{\operatorname{Spec}}
\newcommand\Supp{\operatorname{Supp}}
\newcommand\Proj{\operatorname{Proj}}
\newcommand\Bl{\operatorname{Bl}}
\newcommand\Tot{\operatorname{Tot}}
\newcommand\inv{\operatorname{inv}}
\newcommand\id{\operatorname{id}}
\newcommand\A{\mathsf{A}}
\newcommand\B{\mathsf{B}}
\newcommand\cA{\curly A}
\newcommand{\II}{{I^{\bullet}}}
\newcommand{\JJ}{{J^{\bullet}}}
\newcommand{\Mbar}{{\overline M}}


\newcommand{\Pic}{\mathop{\rm Pic}\nolimits}
\newcommand{\PT}{\mathsf{PT}}
\newcommand{\DT}{\mathsf{DT}}
\newcommand{\GW}{\mathsf{GW}}
\newcommand{\Q}{\mathsf{Q}}
\newcommand{\HH}{\mathrm{H}}
\newcommand{\OO}{\mathscr O}



%%%%%%%%%%%%%%%%%%%%%%%%%% PROOFOF
\makeatletter
\newenvironment{proofof}[1]{\par
  \pushQED{\qed}%
  \normalfont \topsep6\p@\@plus6\p@\relax
  \trivlist
  \item[\hskip3\labelsep
        \itshape
    Proof of #1\@addpunct{.}]\ignorespaces
}{%
  \popQED\endtrivlist\@endpefalse
}
\makeatother


%%%%%%%%%%%%%%%%%%%%%%%%%
\usepackage[all]{xy}
\usepackage{tikz}
\usepackage{tikz-cd}
\usepackage{adjustbox}
\usepackage{rotating}
\usepackage{comment}
\newcommand*{\isoarrow}[1]{\arrow[#1,"\rotatebox{90}{\(\sim\)}"
]}
\usetikzlibrary{matrix,shapes,intersections,arrows,decorations.pathmorphing}
\tikzset{commutative diagrams/arrow style=math font}
\tikzset{commutative diagrams/.cd,
mysymbol/.style={start anchor=center,end anchor=center,draw=none}}
\newcommand\MySymb[2][\square]{%
  \arrow[mysymbol]{#2}[description]{#1}}
\tikzset{
shift up/.style={
to path={([yshift=#1]\tikztostart.east) -- ([yshift=#1]\tikztotarget.west) \tikztonodes}
}
}

\newenvironment{citazione}
  {\begin{quotation}\normalsize}
  {\end{quotation}}
  
%%%%%%%%%%%%%%%%%%%%%%%%%%%%%%%%%%%

\theoremstyle{definition}

\newtheorem*{lemma*}{Lemma}
\newtheorem*{theorem*}{Theorem}
\newtheorem*{example*}{Example}
\newtheorem*{fact*}{Fact}
\newtheorem*{notation*}{Notation}
\newtheorem*{definition*}{Definition}
\newtheorem*{prop*}{Proposition}
\newtheorem*{remark*}{Remark}
\newtheorem*{corollary*}{Corollary}
\newtheorem*{solution}{Solution}
\newtheorem*{conventions*}{Conventions}
\newtheorem*{question*}{Question}
\newtheorem*{convention*}{Convention}

\newtheorem{definition}{Definition}[section]
\newtheorem{assumption}{Assumption}[section]
\newtheorem{problem}[definition]{Problem}
\newtheorem{question}[definition]{Question}
\newtheorem{example}[definition]{Example}
\newtheorem{fact}[definition]{Fact}
%\newtheorem{prop}[definition]{Proposition}
\newtheorem{conventions}[definition]{Conventions}
\newtheorem{convention}{Convention}
\newtheorem{notation}[definition]{Notation}
\newtheorem{remark}[definition]{Remark}
%\newtheorem{theorem}[definition]{Theorem}
%\newtheorem{corollary}[definition]{Corollary}
%\newtheorem{lemma}[definition]{Lemma}
\newtheorem{conjecture}[definition]{Conjecture}
\newtheorem{claim}[definition]{Claim}

\newtheorem{observation}[definition]{Observation}

\newtheoremstyle{thm} % <name> % (ambienti con dimostrazione)
        {3mm}% <Space above>
        {3mm}% <Space below>
        {\slshape}% <Body font> % 
        {0mm}% <Indent amount>
        {\bfseries}% <Theorem head font>
        {.}% <Punctuation after theorem head>
        {1mm}% <Space after theorem head>
        {}% <Theorem head spec (can be left empty, meaning 'normal')> 
\theoremstyle{thm}
\newtheorem{theorem}[definition]{Theorem}
\newtheorem{corollary}[definition]{Corollary}
\newtheorem{lemma}[definition]{Lemma}
\newtheorem{Claim}[definition]{Claim}
\newtheorem{prop}[definition]{Proposition}
\newtheorem{thm}{Theorem}
\renewcommand*{\thethm}{\Alph{thm}}
\newtheorem{conj}{Conjecture}
\renewcommand*{\theconj}{\Alph{conj}}



\newtheoremstyle{ex} % <name> % (ambienti con dimostrazione)
        {3mm}% <Space above>
        {3mm}% <Space below>
        {}% <Body font> % \slshape
        {0mm}% <Indent amount>
        {\scshape}% <Theorem head font>
        {.}% <Punctuation after theorem head>
        {1mm}% <Space after theorem head>
        {}% <Theorem head spec (can be left empty, meaning 'normal')> 
\theoremstyle{ex}
\newtheorem{exercise}[definition]{Exercise}

\newtheoremstyle{sol} % <name> % (ambienti con dimostrazione)
        {3mm}% <Space above>
        {3mm}% <Space below>
        {}% <Body font> % 
        {0mm}% <Indent amount>
        {\scshape}% <Theorem head font>
        {.}% <Punctuation after theorem head>
        {1mm}% <Space after theorem head>
        {}% <Theorem head spec (can be left empty, meaning 'normal')> 
\theoremstyle{sol}









\usepackage{tikz}
\usepackage{xparse}

\newcount\tableauRow
\newcount\tableauCol
\def\tableauDim{0.4}
\newenvironment{Tableau}[1]{%
  \tikzpicture[scale=0.5,draw/.append style={thick,black},
                      baseline=(current bounding box.center)]
    % now draw the tableau
    \tableauRow=-1.5
    \foreach \Row in {#1} {
       \tableauCol=0.5
       \foreach\k in \Row {
         \draw[thin](\the\tableauCol,\the\tableauRow)rectangle++(1,1);
         \draw[thin](\the\tableauCol,\the\tableauRow)+(0.5,0.5)node{$\k$};
         \global\advance\tableauCol by 1
       }
       \global\advance\tableauRow by -1
    }
}{\endtikzpicture}
\newcommand\tableau[1]{\begin{Tableau}{#1}\end{Tableau}}










%Ser
\newtheorem*{Acknowledgments*}{Acknowledgments}